\documentclass[a4paper,12pt]{article}
\usepackage[utf8]{inputenc}
\usepackage{amsmath, amssymb, bm}
\usepackage{geometry}
\geometry{margin=2.5cm}
\usepackage[french]{babel}
\usepackage{enumitem}

\title{Corrections d'exercices de colle -- Mines-Ponts \\ Calcul diff\'erentiel et orthogonalit\'e 2026 La martinière PSI}
\author{}
\date{}

\begin{document}

\maketitle

\section*{Exercice 1 -- Minimum d'une fonction quadratique}

Soit \( A \in \mathcal{S}_n(\mathbb{R}) \) une matrice sym\'etrique \`a spectre strictement positif. Soit \( B \in \mathbb{R}^n \), et on consid\`ere :
\[
f(X) = X^\top A X - 2 B^\top X
\]

\subsection*{1. Calcul du gradient de \( f \)}
Nous voulons calculer le gradient de la fonction \( f : \mathbb{R}^n \to \mathbb{R} \), d\'efinie comme combinaison d'une forme quadratique et d'une forme lin\'eaire. On d\'ecompose :
\[
f(X) = f_1(X) + f_2(X) \quad \text{o\`u} \quad f_1(X) = X^\top A X, \quad f_2(X) = -2 B^\top X
\]

\paragraph{\underline{Étude de \( f_1(X) \)}} Cette fonction est une forme quadratique. Pour calculer son gradient, rappelons que pour toute matrice \( A \in \mathcal{M}_n(\mathbb{R}) \), on a :
\[
\nabla_X (X^\top A X) = (A + A^\top) X
\]
Dans notre cas, \( A \) est sym\'etrique, donc \( A = A^\top \) et la formule devient :
\[
\boxed{\nabla f_1(X) = 2 A X}
\]

\paragraph{Justification explicite par calcul composante par composante :} Soit \( X = (x_1, \dots, x_n)^\top \), alors :
\[
f_1(X) = \sum_{i=1}^n \sum_{j=1}^n x_i A_{ij} x_j
\]
On calcule \( \partial f_1 / \partial x_k \) :
\[
\frac{\partial f_1}{\partial x_k} = \sum_{j=1}^n A_{kj} x_j + \sum_{i=1}^n x_i A_{ik} = \sum_{j=1}^n (A_{kj} + A_{jk}) x_j
\]
Comme \( A \) est sym\'etrique, \( A_{kj} = A_{jk} \), donc :
\[
\frac{\partial f_1}{\partial x_k} = 2 \sum_{j=1}^n A_{kj} x_j = (2AX)_k
\]

\paragraph{\underline{Étude de \( f_2(X) \)}} \( f_2(X) = -2 B^\top X = -2 \sum_{i=1}^n B_i x_i \) est une forme lin\'eaire. On a donc directement :
\[
\boxed{\nabla f_2(X) = -2 B}
\]

\paragraph{\underline{Conclusion :}} En ajoutant les deux parties :
\[
\boxed{\nabla f(X) = 2AX - 2B}
\]

\subsection*{2. Recherche du point critique}
Un point critique de \( f \) est une solution de \( \nabla f(X) = 0 \), soit :
\[
2 A X - 2 B = 0 \quad \Rightarrow \quad A X = B
\]
Comme \( A \succ 0 \), elle est inversible (toutes ses valeurs propres sont strictement positives), donc :
\[
\boxed{X = A^{-1} B}
\]
Ce point est unique et d\'efinit une candidate au minimum global.

\subsection*{3. Valeur du minimum}
On remplace le point critique dans l'expression de \( f \) :
\begin{align*}
f(A^{-1}B) &= (A^{-1}B)^\top A A^{-1}B - 2B^\top A^{-1}B \\
&= B^\top A^{-1} A A^{-1} B - 2 B^\top A^{-1} B \\
&= B^\top A^{-1} B - 2 B^\top A^{-1} B \\
&= \boxed{- B^\top A^{-1} B}
\end{align*}

\subsection*{4. Nature du point critique}
La fonction \( f \) est de classe \( \mathcal{C}^2 \) avec hessienne constante :
\[
\nabla^2 f(X) = 2 A
\]
Or \( A \succ 0 \Rightarrow \nabla^2 f(X) \succ 0 \), donc \( f \) est \textbf{strictement convexe}.
\`A strictement parler, un point critique d'une fonction strictement convexe est un minimum global. Ainsi :
\[
\boxed{X = A^{-1} B \text{ est le minimum global de } f}
\]
\section*{Exercice 2 -- Polynômes orthogonaux dans \(\mathbb{R}_N[X]\)}

On définit :
\[
\varphi_k(X) = \frac{1}{k!} \prod_{i=0}^{k-1}(X - i), \quad
P_n(X) = \sum_{k=0}^n (-1)^k \varphi_k(X) \varphi_{n-k}(N - X)
\]

Et le produit scalaire :
\[
(P|Q) = \sum_{j=0}^N \binom{N}{j} P(j) Q(j)
\]

\subsection*{1. Montrer que \(\varphi_k(n) = \binom{n}{k}\)}

Par définition :
\[
\varphi_k(n) = \frac{1}{k!} \prod_{i=0}^{k-1}(n - i)
= \frac{n(n-1)\cdots(n - k + 1)}{k!} = \binom{n}{k}
\]
On adopte la convention \(\binom{n}{k} = 0\) si \(k > n\), ce qui prolonge la validité du résultat.

\subsection*{2. Montrer que \(f_j(u) = (1 - u)^j (1 + u)^{N - j}\)}

On définit :
\[
f_j(u) = \sum_{n=0}^N P_n(j) u^n
\]
En remplaçant la définition de \(P_n\), on intervertit les sommes :
\[
f_j(u) = \sum_{k=0}^N (-1)^k \varphi_k(j) u^k \cdot \sum_{m=0}^N \varphi_m(N - j) u^m
= \left(\sum_{k=0}^N (-1)^k \binom{j}{k} u^k\right)
  \left(\sum_{m=0}^N \binom{N - j}{m} u^m\right)
\]
\[
= (1 - u)^j (1 + u)^{N - j}
\]

\subsection*{3. Montrer que \(F(u, v) = 2^N (1 + uv)^N\)}

Définition :
\[
F(u, v) = \sum_{j=0}^N \binom{N}{j} f_j(u) f_j(v)
= \sum_{j=0}^N \binom{N}{j} (1 - u)^j (1 + u)^{N - j} (1 - v)^j (1 + v)^{N - j}
\]
\[
= \sum_{j=0}^N \binom{N}{j} [(1 - u)(1 - v)]^j [(1 + u)(1 + v)]^{N - j}
= \left((1 - u)(1 - v) + (1 + u)(1 + v)\right)^N / 2^N
= 2^N (1 + uv)^N
\]

\subsection*{4. Dérivées partielles de \(F\)}

On dérive \(F(u,v) = 2^N (1 + uv)^N\). Alors :
- Si \(a \ne b\), le terme \(\partial_u^a \partial_v^b F(0, 0)\) est nul car il ne peut provenir d’un \((uv)^k\) unique.
- Si \(a = b\), alors :
\[
\partial_u^a \partial_v^a F(0,0) = 2^N \cdot N! \cdot \frac{1}{(N - a)! a! a!} = \frac{2^N N!}{a!(N - a)!}
\]

\subsection*{5. Conclusion : \( (P_k) \) est une base orthogonale}

Comme :
\[
(P_n | P_m) = \partial_u^n \partial_v^m F(0,0)
\Rightarrow
\begin{cases}
0 &\text{si } n \ne m \\\\
\frac{2^N N!}{n!(N - n)!} &\text{si } n = m
\end{cases}
\]
la famille \((P_k)_{0 \leq k \leq N}\) est orthogonale et libre. Elle forme donc une base orthogonale de \(\mathbb{R}_N[X]\).
\section*{Exercice 3 -- Maximum d'une fonction symétrique sur un simplexe}

On travaille sur :
\[
K_n = \left\{ x \in (\mathbb{R}_+)^n \mid \sum x_i = 1 \right\}, \quad
T_n = \left\{ x \in (\mathbb{R}_+^*)^n \mid \sum x_i = 1 \right\}
\]
\[
f_n(x) = \sum_{1 \le i < j \le n} \frac{x_i x_j}{x_i^2 + x_j^2}
\]

\subsection*{1. Justifier l'existence de \(M_n = \max_{x \in K_n} f_n(x)\)}

- \(f_n\) est continue sur \((\mathbb{R}_+^*)^n\), donc sur \(K_n\)
- \(K_n\) est fermé et borné, donc compact
- \(\Rightarrow\) par le théorème de Weierstrass, \(f_n\) atteint son maximum sur \(K_n\)

\subsection*{2. Déterminer \(M_2\)}

\[
f_2(x_1,x_2) = \frac{x_1 x_2}{x_1^2 + x_2^2}
\text{ avec } x_2 = 1 - x_1
\Rightarrow
f_2(x_1) = \frac{x_1(1 - x_1)}{x_1^2 + (1 - x_1)^2}
\]
\[
= \frac{x(1 - x)}{2x^2 - 2x + 1}
\]
On dérive :
\[
f'(x) = \frac{(1 - 2x)(2x^2 - 2x + 1) - x(1 - x)(4x - 2)}{(2x^2 - 2x + 1)^2}
\]
Racine : \(x = \frac{1}{2} \Rightarrow M_2 = f_2\left(\frac{1}{2}\right) = \frac{1/4}{1/2} = \boxed{1/8}\)

\subsection*{3.a Inégalité homogène}

Posons \(y = \frac{x}{\sum x_i} \in T_n\). Alors :
\[
f_n(x) = \left(\sum x_i\right)^4 f_n(y) \Rightarrow
f_n(x) \leq \left(\sum x_i\right)^4 M_n
\]

\subsection*{3.b Gradient et colinéarité}

On maximise \(f_n\) sous contrainte \(\sum x_i = 1\). Le Lagrangien est :
\[
\mathcal{L}(x, \lambda) = f_n(x) - \lambda \left(\sum x_i - 1\right)
\Rightarrow
\nabla f_n(m_0) = \lambda \cdot \nabla \left(\sum x_i\right) = \lambda (1, \dots, 1)
\]

\subsection*{4. Montrer que \(M_n = \frac{1}{8}\)}

Cas optimal : deux coordonnées égales à \(\frac{1}{2}\), autres nulles.
\[
f_n(x) = \frac{\frac{1}{4}}{\frac{1}{2}} = \frac{1}{8}
\Rightarrow \boxed{M_n = \frac{1}{8}} \text{ atteint pour tout } n \geq 2
\]


\vspace{1em}
\hrule
\vspace{1em}

% Les autres exercices suivent après ici...

\end{document}
